\chapter{Complexity — Word}
%%% generated by the blueprint pipeline from TCSlib.Complexity.SpaceComplexity.Codec.Word
\section{Overview}

Declarations of the Lean module \texttt{TCSlib.Complexity.SpaceComplexity.Codec.Word}.

\section{Declarations}

\begin{definition}[Output summary index]
\lean{Complexity.Codec.sumIdx}
\label{Complexity.Codec.sumIdx}
\leanok
\uses{Turing.OutSummary}
Let $N$ be a finite nondeterministic Turing machine with Boolean input symbols. The
empty, accepting, and dead output summaries are assigned indices $0$, $1$, and $2$,
respectively, in $\text{Fin}(3)$.
\end{definition}

\begin{lemma}[Injectivity of the output summary index]
\lean{Complexity.Codec.sumIdx_injective}
\label{Complexity.Codec.sumIdx_injective}
\leanok
\uses{Complexity.Codec.sumIdx}
\difficulty{4}
The map assigning indices $0$, $1$, and $2$ respectively to the empty-output summary,
the summary for the output consisting exactly of $1$, and the summary for all other
outputs is injective.
\end{lemma}

\begin{definition}[Head and output summary pair]
\lean{Complexity.Codec.Hd}
\label{Complexity.Codec.Hd}
\leanok
Let $N$ be a finite nondeterministic Turing machine with Boolean input symbols. A
head-and-summary value consists of either an internal state of $N$ or no state, together
with one of three output-summary indices.
\end{definition}

\begin{definition}[Encoded cell values]
\lean{Complexity.Codec.CellV}
\label{Complexity.Codec.CellV}
\leanok
\uses{Complexity.Codec.Hd}
For a finite nondeterministic Turing machine $N$ with $k$ work tapes, a cell value
consists of an optional head descriptor and, for each work tape, a pair comprising an
optional Boolean symbol and a Boolean marker.
\end{definition}

\begin{definition}[Number of head descriptors]
\lean{Complexity.Codec.hH}
\label{Complexity.Codec.hH}
\leanok
\uses{Complexity.Codec.Hd}
For a finite nondeterministic Turing machine $N$, let $h$ be the number of possible head
descriptors of $N$.
\end{definition}

\begin{definition}[Encoding chunk width]
\lean{Complexity.Codec.W}
\label{Complexity.Codec.W}
\leanok
\uses{Complexity.Codec.hH}
For a finite nondeterministic Turing machine $N$ with $k$ work tapes, define the chunk
width $W$ by $W=h+3k+1$, where $h$ is the number of possible head descriptors.
\end{definition}

\begin{definition}[Head descriptor index]
\lean{Complexity.Codec.hdNum}
\label{Complexity.Codec.hdNum}
\leanok
\uses{Complexity.Codec.Hd}
For a finite nondeterministic Turing machine $N$, a head descriptor is a pair consisting
of either a machine state or no state, together with one of three indices. Assign each
such descriptor its index in a fixed enumeration of all head descriptors, starting at
zero.
\end{definition}

\begin{lemma}[Bound on Head-Descriptor Numbers]
\lean{Complexity.Codec.hdNum_lt}
\label{Complexity.Codec.hdNum_lt}
\leanok
\uses{Complexity.Codec.Hd, Complexity.Codec.hH, Complexity.Codec.hdNum}
\difficulty{1}
Let $N$ be a finite nondeterministic Turing machine with Boolean input alphabet, and let
$h$ be a head descriptor, consisting of an optional state and one of three output-status
values. Its number in the fixed enumeration of head descriptors is strictly less than
the total number of such descriptors.
\end{lemma}

\begin{lemma}[Injective Numbering of Head Descriptors]
\lean{Complexity.Codec.hdNum_injective}
\label{Complexity.Codec.hdNum_injective}
\leanok
\uses{Complexity.Codec.Hd, Complexity.Codec.hdNum}
\difficulty{3}
Let $N$ be a finite nondeterministic Turing machine with Boolean input alphabet. A head
descriptor consists of an optional machine state together with one of three
output-status values; assigning to each descriptor its number in a fixed enumeration of
all such descriptors is injective.
\end{lemma}

\begin{definition}[Bit encoding of a cell]
\lean{Complexity.Codec.chunkBit}
\label{Complexity.Codec.chunkBit}
\leanok
\uses{Complexity.Codec.CellV, Complexity.Codec.hH, Complexity.Codec.hdNum}
Let $N$ be a finite nondeterministic Turing machine with $k$ work tapes, and let a cell
value contain an optional head descriptor and, for each tape, an optional Boolean symbol
and a Boolean marker. For each natural number $t$, define its encoded bit as follows: if
$t<h$, it records whether the head descriptor's index is $t$; otherwise, for tape index
$i=rac{t-h}{3}$ when $i<k$, the remainder of $t-h$ upon division by $3$ selects,
respectively, whether that tape's symbol is present, its symbol with absence interpreted
as false, or its marker. If the resulting tape index is at least $k$, the bit is true.
\end{definition}

\begin{definition}[Encoding one cell as a bit string]
\lean{Complexity.Codec.encChunk}
\label{Complexity.Codec.encChunk}
\leanok
\uses{Complexity.Codec.CellV, Complexity.Codec.W, Complexity.Codec.chunkBit}
Let $N$ be a finite nondeterministic Turing machine, and let $c$ be a cell value. Encode
$c$ as the Boolean string of length $W$, whose bit at each position is the bit encoding
of $c$ at that position.
\end{definition}

\begin{lemma}[Length of a cell encoding]
\lean{Complexity.Codec.length_encChunk}
\label{Complexity.Codec.length_encChunk}
\leanok
\uses{Complexity.Codec.CellV, Complexity.Codec.W, Complexity.Codec.encChunk}
\difficulty{1}
For a finite-state nondeterministic Turing machine $N$ and any cell value $c$, the
encoding of $c$ has length exactly $W$, the cell-code width for $N$.
\end{lemma}

\begin{lemma}[Encoded cell bit retrieval]
\lean{Complexity.Codec.getD_encChunk}
\label{Complexity.Codec.getD_encChunk}
\leanok
\uses{Complexity.Codec.CellV, Complexity.Codec.W, Complexity.Codec.chunkBit, Complexity.Codec.encChunk}
\difficulty{1}
Let $N$ be a finite nondeterministic Turing machine with Boolean input symbols, and let
$cv$ be a cell value for $N$. For every natural number $t$ with $t<W(N)$, the bit at
position $t$ of the length-$W(N)$ encoding of $cv$, using \text{false} as the default for
an out-of-range list lookup, equals the bit specified by the chunk-bit rule at position
$t$.
\end{lemma}

\begin{lemma}[State-slot bit in a chunk]
\lean{Complexity.Codec.chunkBit_slot}
\label{Complexity.Codec.chunkBit_slot}
\leanok
\uses{Complexity.Codec.CellV, Complexity.Codec.chunkBit, Complexity.Codec.hH, Complexity.Codec.hdNum}
\difficulty{1}
Let $N$ be a finite nondeterministic Turing machine with Boolean tape symbols, let $v$
be a cell value, and let $t$ be a natural number less than the number of possible
state-summary pairs. The bit of the chunk encoding of $v$ at position $t$ is true
exactly when the optional state-summary pair recorded in $v$ has numerical code $t$.
\end{lemma}

\begin{lemma}[Work-tape slots in a chunk]
\lean{Complexity.Codec.chunkBit_tape}
\label{Complexity.Codec.chunkBit_tape}
\leanok
\uses{Complexity.Codec.CellV, Complexity.Codec.chunkBit, Complexity.Codec.hH}
\difficulty{4}
Let $N$ be a finite nondeterministic Turing machine with Boolean tape symbols, let $v$
be a cell value, let $i$ be a work-tape index, and let $r\in\mathbb{N}$ satisfy $r<3$.
The chunk bit at position $h_H+3i+r$ records, according as $r=0$, $r=1$, or otherwise,
whether the tape cell is nonblank, its Boolean symbol (taking false when no symbol is
present), or whether the work-tape head is at that cell.
\end{lemma}

\begin{theorem}[Injectivity of chunk encoding]
\lean{Complexity.Codec.encChunk_injective}
\label{Complexity.Codec.encChunk_injective}
\leanok
\uses{Complexity.Codec.Hd, Complexity.Codec.W, Complexity.Codec.chunkBit, Complexity.Codec.chunkBit_slot, Complexity.Codec.chunkBit_tape, Complexity.Codec.encChunk, Complexity.Codec.getD_encChunk, Complexity.Codec.hH, Complexity.Codec.hdNum, Complexity.Codec.hdNum_injective, Complexity.Codec.hdNum_lt}
\difficulty{6}
For every finite nondeterministic Turing machine $N$, the encoding of cell values as
chunks is injective: if two cell values have the same encoded chunk, then the cell
values are equal.
\end{theorem}

\begin{lemma}[Extracting a block from a flattened list]
\lean{Complexity.Codec.take_drop_flatten}
\label{Complexity.Codec.take_drop_flatten}
\leanok
\difficulty{4}
Let $L$ be a finite list of lists, each of length $w$, and let $j$ be an index with
$j<|L|$. In the concatenation of all lists in $L$, the block of length $w$ beginning at
position $wj$ is exactly the $j$th list.
\end{lemma}

\begin{lemma}[Length of a concatenation of equal-length words]
\lean{Complexity.Codec.length_flatten_const}
\label{Complexity.Codec.length_flatten_const}
\leanok
\difficulty{3}
Let $L$ be a finite list of words over an alphabet, and let $w$ be a natural number. If
every word in $L$ has length $w$, then the concatenation of the words in $L$ has length
$w|L|$.
\end{lemma}

\begin{definition}[Encoding a sequence of cells]
\lean{Complexity.Codec.chunksWord}
\label{Complexity.Codec.chunksWord}
\leanok
\uses{Complexity.Codec.CellV, Complexity.Codec.encChunk}
Let $N$ be a finite nondeterministic Turing machine, let $f$ assign a cell value to each
natural number, and let $m$ be a natural number. Encode the first $m$ values of $f$ by
concatenating their cell encodings in order.
\end{definition}

\begin{lemma}[Length of a word of encoded cells]
\lean{Complexity.Codec.length_chunksWord}
\label{Complexity.Codec.length_chunksWord}
\leanok
\uses{Complexity.Codec.CellV, Complexity.Codec.W, Complexity.Codec.chunksWord, Complexity.Codec.length_flatten_const}
\difficulty{3}
Let $N$ be a finite-state nondeterministic Turing machine, let $f$ assign a cell value
to each natural number, and let $m$ be a natural number. The word formed by
concatenating the encodings of $f(0),\ldots,f(m-1)$ has length $Wm$, where $W$ is the
cell-code width for $N$.
\end{lemma}

\begin{lemma}[Injectivity of concatenated chunk encodings]
\lean{Complexity.Codec.chunksWord_inj}
\label{Complexity.Codec.chunksWord_inj}
\leanok
\uses{Complexity.Codec.CellV, Complexity.Codec.W, Complexity.Codec.chunksWord, Complexity.Codec.encChunk, Complexity.Codec.encChunk_injective, Complexity.Codec.take_drop_flatten}
\difficulty{4}
Let $N$ be a finite nondeterministic Turing machine with Boolean tape symbols, and let
$f,g$ assign a cell value to each natural number. If the words formed by concatenating
the first $m$ encoded cell values of $f$ and $g$ are equal, then $f(j)=g(j)$ for every
$j<m$.
\end{lemma}

\begin{definition}[Cell value from a configuration]
\lean{Complexity.Codec.cellOf}
\label{Complexity.Codec.cellOf}
\leanok
\uses{Complexity.Codec.CellV, Complexity.Codec.sumIdx, Turing.Cfg, Turing.outSummary}
Let $N$ be a finite nondeterministic Turing machine with $k$ work tapes, let $x$ be a
Boolean input string, let $C$ be a configuration of $N$ on input $x$, and let $s,j$ be
natural numbers. The cell value at position $j$ has an optional head descriptor: it is
present when $j=0$ and then records the machine state together with the three-valued
summary of the output tape; otherwise it is absent. For each work tape, it records the
symbol at position $j-s$ and whether that tape's head is at position $j-s$.
\end{definition}

\begin{definition}[Binary width of a natural number]
\lean{Complexity.Codec.wP}
\label{Complexity.Codec.wP}
\leanok
For each natural number $n$, let $w(n)$ be the number of binary digits in the usual
binary expansion of $n+1$.
\end{definition}

\begin{definition}[Right-padded binary representation]
\lean{Complexity.Codec.padBits}
\label{Complexity.Codec.padBits}
\leanok
For natural numbers $w$ and $v$, form the binary digit list of $v$ and append enough
false bits to make its length $w$ when the digit list has length at most $w$. If it
already has length greater than $w$, append no bits.
\end{definition}

\begin{definition}[Configuration encoding word]
\lean{Complexity.Codec.cfgWord}
\label{Complexity.Codec.cfgWord}
\leanok
\uses{Complexity.Codec.cellOf, Complexity.Codec.chunksWord, Complexity.Codec.padBits, Complexity.Codec.wP, Turing.Cfg}
Let $N$ be a finite nondeterministic Turing machine with Boolean tape symbols, let $x$
be a Boolean input string, and let $C$ be a configuration of $N$ on input $x$. For a
natural number $s$, the encoding word is the concatenation of the encoded tape cells in
the window determined by $s$, followed by the input-head position written as a binary
list padded to width $w(|x|)$, and then a final true bit.
\end{definition}

\begin{definition}[Natural-number configuration code]
\lean{Complexity.Codec.cfgCode}
\label{Complexity.Codec.cfgCode}
\leanok
\uses{Complexity.Codec.cfgWord, Complexity.LogProg.bitsVal, Turing.Cfg}
Let $N$ be a finite nondeterministic Turing machine with Boolean tape symbols, let $x$
be a Boolean input string, and let $C$ be a configuration of $N$ on input $x$. For a
natural number $s$, its code is the natural number represented by the configuration
encoding word at window parameter $s$.
\end{definition}

\begin{lemma}[Appending False Bits Preserves Value]
\lean{Complexity.Codec.bitsVal_append_replicate}
\label{Complexity.Codec.bitsVal_append_replicate}
\leanok
\uses{Complexity.LogProg.bitsVal}
\difficulty{4}
For every finite Boolean list $w$ and every natural number $z$, appending $z$ false
entries to $w$ leaves its value under the binary value function unchanged.
\end{lemma}

\begin{lemma}[Binary expansion with zero padding]
\lean{Complexity.Codec.eq_bits_append_replicate}
\label{Complexity.Codec.eq_bits_append_replicate}
\leanok
\uses{Complexity.LogProg.bitsVal, Complexity.LogProg.bits_two_mul_add}
\difficulty{6}
For every finite Boolean list $w$, the list is the binary expansion of its value, with
enough trailing false entries appended to make its length equal to the length of $w$.
Here the list represents binary digits in least-significant-first order.
\end{lemma}

\begin{lemma}[Length of a padded bit string]
\lean{Complexity.Codec.length_padBits}
\label{Complexity.Codec.length_padBits}
\leanok
\uses{Complexity.Codec.padBits}
\difficulty{2}
For natural numbers $w$ and $v$, if the binary expansion of $v$ has length at most $w$,
then the bit string obtained by appending enough zeros to that expansion to reach length
$w$ has length exactly $w$.
\end{lemma}

\begin{lemma}[Value of a padded binary expansion]
\lean{Complexity.Codec.bitsVal_padBits}
\label{Complexity.Codec.bitsVal_padBits}
\leanok
\uses{Complexity.Codec.bitsVal_append_replicate, Complexity.Codec.padBits, Complexity.LogProg.bitsVal, Complexity.LogProg.bitsVal_bits}
\difficulty{1}
For all natural numbers $w$ and $v$, the binary value of the binary expansion of $v$
followed by enough false bits to make the prescribed padding length $w$ is exactly $v$.
\end{lemma}

\begin{lemma}[List equals its padded binary representation]
\lean{Complexity.Codec.eq_padBits}
\label{Complexity.Codec.eq_padBits}
\leanok
\uses{Complexity.Codec.eq_bits_append_replicate, Complexity.Codec.padBits, Complexity.LogProg.bitsVal}
\difficulty{3}
Let $p$ be a finite Boolean list. Then $p$ equals the least-significant-first binary
expansion of its value, followed by false entries until the resulting list has length
$|p|$.
\end{lemma}

\begin{lemma}[Bound on binary representation length]
\lean{Complexity.Codec.length_bits_pos_le}
\label{Complexity.Codec.length_bits_pos_le}
\leanok
\uses{Complexity.Codec.wP, Complexity.LogProg.length_bits_mono}
\difficulty{1}
For natural numbers $v$ and $n$, if $v\le n+1$, then the length of the binary word
representing $v$ is at most the length of the binary word representing $n+1$.
\end{lemma}

\begin{lemma}[Length of the configuration encoding]
\lean{Complexity.Codec.length_cfgWord}
\label{Complexity.Codec.length_cfgWord}
\leanok
\uses{Complexity.Codec.W, Complexity.Codec.cfgWord, Complexity.Codec.length_bits_pos_le, Complexity.Codec.length_chunksWord, Complexity.Codec.length_padBits, Complexity.Codec.wP, Turing.Cfg}
\difficulty{3}
Let $N$ be a finite-state nondeterministic Turing machine with $k$ work tapes, let $x$
be a Boolean input word, let $s$ be a natural number, and let $C$ be a configuration of
$N$ on input $x$. The encoded configuration word has length $W(2s+1)+w_P(|x|)+1$, where
$W$ is the fixed cell-code width for $N$ and $w_P(n)$ is the length of the binary
representation of $n+1$.
\end{lemma}

\begin{lemma}[Canonical configuration encoding]
\lean{Complexity.Codec.canon_cfgWord}
\label{Complexity.Codec.canon_cfgWord}
\leanok
\uses{Complexity.Codec.cfgWord, Complexity.LogProg.Canon, Turing.Cfg}
\difficulty{1}
Let $N$ be a finite nondeterministic Turing machine with Boolean symbols, let $x$ be a
finite Boolean string, let $s$ be a natural number, and let $C$ be a configuration of
$N$ on input $x$. The Boolean string encoding $C$ in a window of radius $s$ is nonempty
and has final bit true.
\end{lemma}

\begin{theorem}[Binary expansion of a configuration code]
\lean{Complexity.Codec.bits_cfgCode}
\label{Complexity.Codec.bits_cfgCode}
\leanok
\uses{Complexity.Codec.canon_cfgWord, Complexity.Codec.cfgCode, Complexity.Codec.cfgWord, Complexity.LogProg.canon_eq_bits, Turing.Cfg}
\difficulty{1}
Let $N$ be a finite nondeterministic Turing machine over Boolean symbols, let $x$ be a
finite Boolean input, let $s$ be a natural number, and let $C$ be a configuration of $N$
on input $x$. The binary expansion of the numerical code of $C$ with space parameter $s$
is exactly the Boolean word encoding that configuration with parameter $s$.
\end{theorem}

\begin{lemma}[Cell encoding depends on core data]
\lean{Complexity.Codec.cellOf_congr}
\label{Complexity.Codec.cellOf_congr}
\leanok
\uses{Complexity.Codec.cellOf, Turing.Cfg, Turing.MultiTapeTM.ConfigCount.core, Turing.NDTM.coreSum}
\difficulty{4}
Let $N$ be a finite nondeterministic Turing machine with Boolean symbols, let $x$ be a
finite Boolean string, let $s$ be a natural number, and let $C,D$ be configurations on
input $x$. If $C$ and $D$ agree in their state, input-head position, work-tape contents,
work-tape head positions, and output summary, then their encoded cell functions on the
window of radius $s$ are equal.
\end{lemma}

\begin{theorem}[Configuration word congruence]
\lean{Complexity.Codec.cfgWord_congr}
\label{Complexity.Codec.cfgWord_congr}
\leanok
\uses{Complexity.Codec.cellOf_congr, Complexity.Codec.cfgWord, Turing.Cfg, Turing.MultiTapeTM.ConfigCount.core, Turing.NDTM.coreSum}
\difficulty{3}
Let $N$ be a finite nondeterministic Turing machine with Boolean tape symbols, let $x$
be a Boolean input word, and let $C,D$ be configurations of $N$ on input $x$. For every
$s\in\mathbb{N}$, if the configurations have the same core data and the same summary of
their output words, then their encoded configuration words are equal.
\end{theorem}

\begin{theorem}[Recovery of the configuration core]
\lean{Complexity.Codec.coreSum_eq_of_cfgWord_eq}
\label{Complexity.Codec.coreSum_eq_of_cfgWord_eq}
\leanok
\uses{Complexity.Codec.bitsVal_padBits, Complexity.Codec.cellOf, Complexity.Codec.cfgWord, Complexity.Codec.chunksWord, Complexity.Codec.chunksWord_inj, Complexity.Codec.length_chunksWord, Complexity.Codec.sumIdx_injective, Complexity.LogProg.bitsVal, Turing.Cfg, Turing.Cfg.InWindow, Turing.MultiTapeTM.ConfigCount.core, Turing.NDTM.coreSum}
\difficulty{6}
Let $N$ be a finite nondeterministic Turing machine with Boolean tape symbols, let $x$
be a Boolean input, and let $s$ be a natural number. Suppose configurations $C$ and $D$
have every work-tape head within distance $s$ of the origin, and every work-tape cell
outside that interval is blank. If their configuration words are equal, then their cores
together with their output summaries are equal.
\end{theorem}

\begin{definition}[Machine-dependent code length]
\lean{Complexity.Codec.codeK}
\label{Complexity.Codec.codeK}
\leanok
\uses{Complexity.Codec.W, Turing.FinNDTM}
For a finite nondeterministic Turing machine $N$ with Boolean tape symbols, define its
code-length parameter to be $2W(N)+2$, where $W(N)$ is the machine quantity denoted by
$W$ in the surrounding definitions.
\end{definition}

\begin{theorem}[Configuration encoding length bound]
\lean{Complexity.Codec.length_cfgWord_le}
\label{Complexity.Codec.length_cfgWord_le}
\leanok
\uses{Complexity.Codec.W, Complexity.Codec.cfgWord, Complexity.Codec.codeK, Complexity.Codec.length_cfgWord, Complexity.Codec.wP, Complexity.LogProg.length_bits_le_log, Complexity.logSpace, Turing.Cfg}
\difficulty{5}
Let $N$ be a finite nondeterministic Turing machine over the Boolean alphabet, with $k$
work tapes, and let $x$ be a finite Boolean word. For every natural number $s$ and every
configuration $C$ of $N$ on input $x$, the encoding of $C$ has length at most
$(2W+2)(s+\\text{logSpace}(|x|)+1)$, where $W$ is the number of bits used to encode one
tape cell and $\text{logSpace}(n)=\lfloor\log_2 n\rfloor+1$.
\end{theorem}

\begin{lemma}[Value bounded by the corresponding power of two]
\lean{Complexity.Codec.bitsVal_lt_two_pow}
\label{Complexity.Codec.bitsVal_lt_two_pow}
\leanok
\uses{Complexity.LogProg.bitsVal}
\difficulty{4}
For every finite Boolean list $w$, its binary value is strictly less than $2^{|w|}$,
where $|w|$ is the length of the list.
\end{lemma}

\begin{theorem}[Configuration code size bound]
\lean{Complexity.Codec.cfgCode_lt}
\label{Complexity.Codec.cfgCode_lt}
\leanok
\uses{Complexity.Codec.bitsVal_lt_two_pow, Complexity.Codec.cfgCode, Complexity.Codec.codeK, Complexity.Codec.length_cfgWord_le, Complexity.logSpace, Turing.Cfg}
\difficulty{1}
Let $N$ be a finite nondeterministic Turing machine with Boolean symbols, let $x$ be a
finite Boolean string, let $s$ be a natural number, and let $C$ be a configuration of
$N$ on input $x$. Let $K=2W+2$, where $W$ is the number of bits allotted to encode one
cell, and let $L(n)=\lfloor\log_2 n\rfloor+1$ for positive $n$ and $L(0)=1$. The natural
number represented by the Boolean encoding of $C$ in a window of radius $s$ is strictly
less than $2^{K(s+L(|x|)+1)}$.
\end{theorem}

\begin{definition}[Configuration decoder]
\lean{Complexity.Codec.decodeCfg}
\label{Complexity.Codec.decodeCfg}
\leanok
\uses{Complexity.Codec.cfgCode, Turing.Cfg, Turing.Cfg.InWindow}
Let $N$ be a finite nondeterministic Turing machine with Boolean tape symbols, and let
$x$ be a Boolean input string. Given natural numbers $s$ and $c$, return a configuration
of $N$ on input $x$ whose configuration code at window parameter $s$ is $c$ and which
lies in the corresponding configuration window, if such a configuration exists;
otherwise return no configuration.
\end{definition}

\begin{theorem}[Configuration code decoding]
\lean{Complexity.Codec.decodeCfg_cfgCode}
\label{Complexity.Codec.decodeCfg_cfgCode}
\leanok
\uses{Complexity.Codec.bits_cfgCode, Complexity.Codec.cfgCode, Complexity.Codec.coreSum_eq_of_cfgWord_eq, Complexity.Codec.decodeCfg, Turing.Cfg, Turing.Cfg.InWindow, Turing.NDTM.coreSum}
\difficulty{3}
Let $N$ be a finite nondeterministic Turing machine, let $x$ be a Boolean input, let $s$
be a natural number, and let $C$ be a configuration whose work-tape heads lie within
distance $s$ of the origin and whose work tapes are blank outside that window. Decoding
the code of $C$ at window size $s$ returns some configuration $D$ that also lies within
the window, has the same code as $C$, and has the same state, input-head position,
work-tape contents and head positions, and output summary as $C$.
\end{theorem}

\begin{theorem}[Characterization of failed configuration decoding]
\lean{Complexity.Codec.decodeCfg_eq_none}
\label{Complexity.Codec.decodeCfg_eq_none}
\leanok
\uses{Complexity.Codec.cfgCode, Complexity.Codec.decodeCfg, Turing.Cfg, Turing.Cfg.InWindow}
\difficulty{1}
For a finite nondeterministic Turing machine $N$, Boolean input $x$, and natural numbers
$s,c$, decoding $c$ at window size $s$ returns no configuration if and only if there is
no configuration for input $x$ whose work-tape heads lie within distance $s$ of the
origin, whose work tapes are blank outside that window, and whose code at window size
$s$ is $c$.
\end{theorem}
